\documentclass[11pt, a4paper]{report}

\usepackage[utf8]{inputenc}
\usepackage[T1]{fontenc}
\usepackage[french]{babel}

% Appel des packages dans le dossier preambule
\def\packagepath{../../../../preambule}   % path du package princial

\usepackage{\packagepath/page_layout} % <--- Nouveau
\usepackage{\packagepath/config_info}
\usepackage{\packagepath/cadres_cours}
\usepackage{\packagepath/zones_reponse}
\usepackage{\packagepath/config_ds}
\usepackage{\packagepath/config_maths}

% --- PILOTAGE ---
\def\visibleornot{invisible}    % visible or invisible


\def\documentpath{./Sources_Latex}
\graphicspath{{./Images/}}


\def\ptsexoA{0}

\begin{document}

%%%%%%%%%%%%%%%%%%%%%%%%%%%
\def\level{BTS}  % Classe à droite
\def\course{CIEL 1}
\def\chaptername{Intégration} % Titre au centre
\def\docname{C107 -- TP2}        % Identifiant à gauche

\pagestyle{DS_FP}
%%%%%%%%%%%%%%%%%%%%%%%%%%%%%%%%%

\NomPrenomNote{}

%%=============================================================

\section*{PARTIE A: Détermination expérimentale du paramètre $\tau$}

On étudie la tension de décharge d'un composant modélisée par la fonction $f$ définie sur $[0; +\infty[$ par la relation suivante:
\[f(x) = 10\cdot \e^{-x/\tau} \qquad \text{ (où $\tau$ est la constante de temps).}\]

Le signal doit respecter les contraintes physiques suivantes:

\begin{itemize}
    \item À l'instant $x=0$, la tension doit être de 10V, soit $f(0) = 10$.
    \item À l'instant $x=2$, la pente de la tangente à la courbe doit être égale à $-1,84$.
\end{itemize}

\begin{enumerate}
    \item Tracer la courbe de $f$ pour différentes valeurs de $\tau$.  
    \item Identifier la valeur de $\tau$ qui satisfait les contraintes données.  
    \begin{blocvisible}[height=1]
        $\tau = 2$ satisfait les contraintes, car $f(0) = 10$ et $f'(2) = -1,84$.
    \end{blocvisible}
\end{enumerate}


\section*{PARTIE B: Étude du signal identifié}

Utilisez maintenant la valeur trouvée ($\tau = 2$) pour analyser le comportement du signal $f(x) = 10\cdot \e^{-0,5x}$.

\begin{enumerate}[resume]
    \item Tracer la courbe de $f$ et observer son comportement.  
    \item Affichez la fonction dérivée $f'$   
    \item Déterminer le signe de $f'$ et en déduire les variations de $f$.  
    \begin{blocvisible}[height=4]
        \begin{minipage}{0.54\linewidth}
            $f'(x) = -5\cdot \e^{-0,5x}$, qui est toujours négatif pour $x \geq 0$.

            Par conséquent, $f$ est strictement décroissante sur $\fo{0}{+\infty}$.
        \end{minipage}\hfill
        \begin{minipage}{0.45\linewidth}
            \begin{tikzpicture}
                \tkzTabInit[lgt=2.5,espcl=3.5]
                {$x$ / 1, $signe \, f'(x)$ / 1, $var \, f$ / 1}
                {$0$, $+\infty$}
                \tkzTabLine{}
                \tkzTabVar{}
            \end{tikzpicture}
        \end{minipage}
    \end{blocvisible}
        
    \item Placer un point $N\left(x_N; f(x_N)\right)$ sur la courbe et observer la valeur de $f(x_N)$ pour différentes positions de $N$.   
    \item Déplacez le point $N$ vers la droite (grandes valeurs de $x$) et observez la tendance de $f(x_N)$.  
    \item Conjecturez la limite de $f(x)$ lorsque $x$ tend vers $+\infty$.   
    
    \begin{blocvisible}[height=2]
        Lorsque $x$ tend vers $+\infty$, $f(x)$ tend vers 0, car la fonction exponentielle décroît rapidement.

        Cela signifie que la tension de décharge approche de 0V à mesure que le temps passe, ce qui est cohérent avec le comportement attendu d'un signal de décharge.
    \end{blocvisible}
\end{enumerate}

\section*{PARTIE C: Intégration et Valeur Moyenne}

On souhaite évaluer la tension moyenne délivrée par le composant sur un intervalle $\ff{a}{b}$.

\begin{enumerate}[resume]
    \item Créer les curseurs $a$ et $b$ avec les valeurs initiales $a=0$ et $b=4$.   
    \item Déterminer graphiquement l'aire sous la courbe de $f$ entre $a$ et $b$ avec la commande \verb|Aire = Intégrale(f, a, b)|.   
    \begin{blocvisible}[height=2]
        D'après Géogébra, l'aire sous la courbe de $f$ entre $a=0$ et $b=4$ est d'environ $13,53$ unités d'aires.
    \end{blocvisible}       
    
    \pagebreak
    \thispagestyle{DS_LP}

    \item Observer comment l'aire change lorsque vous modifiez les valeurs de $a$ et $b$.  
    \begin{blocvisible}[height=2]
        Lorsque vous augmentez $b$, l'aire sous la courbe de $f$ entre $a$ et $b$ augmente, car vous incluez plus de la courbe. Inversement, lorsque vous augmentez $a$, l'aire diminue, car vous excluez une partie de la courbe.
    \end{blocvisible}
    \item Calculer l'aire sous la courbe de $f$ entre $a=0$ et $b=4$ à l'aide de l'intégrale.  
    \begin{blocvisible}[height=3]
        \[A=\int_a^b f(x) \dx = \int_0^4 10\cdot \e^{-0,5x} \dx = \left[-20\cdot \e^{-0,5x} \right]_0^4 =  20 \cdot (\e^{-0,5\times 4} - \e^{-0,5\times 0}) = 20 \cdot (1 - \e^{-2}) = 13,53\]

        L'aire calculée à l'aide de l'intégrale est d'environ $13,53$ unités d'aires, ce qui correspond à la valeur obtenue graphiquement.
    \end{blocvisible}
    \item Déterminer la valeur moyenne $\mu$ de la fonction $f$ sur l'intervalle $[a ; b]$.  
    \begin{blocvisible}[height=2]           
        \[\mu = \frac{1}{b-a} \int_a^b f(x) \dx = \frac{1}{4-0} \cdot 13,53 = 3,38\]

        La valeur moyenne de la fonction $f$ sur l'intervalle $\ff{0}{4}$ est d'environ $3,38$ volts.
    \end{blocvisible}
    \item Créer un rectangle d'équivalence avec les points \verb|(a,0)|, \verb|(b,0)|,\verb|(b, moy)| et \verb|(a, moy)|. 
    
    Utilisez l'outil \verb|Polygone| pour le colorer.  
\end{enumerate}


\end{document}